\section{Mixed Signal Simulation in
ngspice}\label{mixed-signal-simulation-in-ngspice}

\textbf{Keywords:} Mixed Signal, ngspice, Verilog, RTL, d\_cosim,
Testbench

\section{Digital simulation}\label{digital-simulation}\index{digital simulation@Digital simulation}

\begin{itemize}
\item
  The order of execution of events at the same time-step does not matter
\item
  The system is causal. Changes in the future do not affect signals in
  the past or the now
\end{itemize}

There are both commercial and open source tools for digital simulation.
If you've never used a digital simulator, then I'd recommend you start
with iverilog. I've made some examples at
\href{https://github.com/wulffern/dicex/tree/main/project/verilog}{dicex}.

\textbf{Commercial}

\begin{itemize}
\tightlist
\item
  \href{https://www.cadence.com/ko_KR/home/tools/system-design-and-verification/simulation-and-testbench-verification/xcelium-simulator.html}{Cadence
  Excelium}
\item
  \href{https://eda.sw.siemens.com/en-US/ic/questa/simulation/advanced-simulator/}{Siemens
  Questa}
\item
  \href{https://www.synopsys.com/verification/simulation/vcs.html}{Synopsys
  VCS}
\end{itemize}

\textbf{Open Source}

\begin{itemize}
\tightlist
\item
  \href{https://github.com/steveicarus/iverilog}{iverilog/vpp}
\item
  \href{https://www.veripool.org/verilator/}{Verilator}
\item
  \href{https://sourceforge.net/projects/systemdotnet/}{SystemDotNet}
\end{itemize}

Below is an example of a counter in SystemVerilog. The code can be found
at
\href{https://github.com/wulffern/dicex/tree/main/sim/verilog/counter_sv}{counter\_sv}.

In the always\_ff sections we code what will become our registers;
combinatorial logic would go in an always\_comb section.

\begin{Shaded}
\begin{Highlighting}[]
\KeywordTok{module}\NormalTok{ dig}\OperatorTok{(}
           \DataTypeTok{input} \DataTypeTok{wire}\NormalTok{         clk}\OperatorTok{,}
           \DataTypeTok{input} \DataTypeTok{wire}\NormalTok{         reset}\OperatorTok{,}
           \DataTypeTok{output}\NormalTok{ logic }\OperatorTok{[}\DecValTok{4}\OperatorTok{:}\DecValTok{0}\OperatorTok{]}\NormalTok{ b}
           \OperatorTok{);}

\NormalTok{   logic                      rst }\OperatorTok{=} \DecValTok{0}\OperatorTok{;}

\NormalTok{   always\_ff }\OperatorTok{@(}\KeywordTok{posedge}\NormalTok{ clk}\OperatorTok{)} \KeywordTok{begin}
      \KeywordTok{if}\OperatorTok{(}\NormalTok{reset}\OperatorTok{)}
\NormalTok{        rst }\OperatorTok{\textless{}=} \DecValTok{1}\OperatorTok{;}
      \KeywordTok{else}
\NormalTok{        rst }\OperatorTok{\textless{}=} \DecValTok{0}\OperatorTok{;}
   \KeywordTok{end}

\NormalTok{   always\_ff }\OperatorTok{@(}\KeywordTok{posedge}\NormalTok{ clk}\OperatorTok{)} \KeywordTok{begin}
      \KeywordTok{if}\OperatorTok{(}\NormalTok{rst}\OperatorTok{)}
\NormalTok{        b }\OperatorTok{\textless{}=} \DecValTok{0}\OperatorTok{;}
      \KeywordTok{else}
\NormalTok{        b }\OperatorTok{\textless{}=}\NormalTok{ b }\OperatorTok{+} \DecValTok{1}\OperatorTok{;}
   \KeywordTok{end} \CommentTok{// dig}

\KeywordTok{endmodule}
\end{Highlighting}
\end{Shaded}

\section{Transient analog simulation}\label{transient-analog-simulation}\index{transient analog simulation@Transient analog simulation}

Analog simulation is different. There is no quantized time step. How
fast ``things'' happen in the circuit is entirely determined by the time
constants, change in voltage, and change in current in the system.

It is possible to have a fixed time-step in analog simulation, for
example, we say that nothing is faster than 1 fs, so we pick that as our
time step. If we wanted to simulate 1 s, however, that's at least 1e15
events, and with 1 event per microsecond on a computer it's still a
simulation time of 31 years. Not a viable solution for all analog
circuits.

Analog circuits are also non-linear, properties of resistors,
capacitors, inductors, diodes may depend on the voltage or current
across, or in, the device. Solving for all the non-linear differential
equations is tricky.

An analog simulation engine must parse spice netlist, and setup
partial/ordinary differential equations for node matrix

The nodal matrix could look like the matrix below, \(i\) are the
currents, \(v\) the voltages, and \(G\) the conductances between nodes.

\[
\begin{pmatrix}
G_{11} &G_{12}  &\cdots  &G_{1N} \\ 
G_{21} &G_{22}  &\cdots  &G_{2N} \\ 
\vdots &\vdots  &\ddots  & \vdots\\ 
G_{N1} &G_{N2}  &\cdots  &G_{NN} 
\end{pmatrix}
\begin{pmatrix}
v_1\\ 
v_2\\ 
\vdots\\ 
v_N
\end{pmatrix}=
\begin{pmatrix}
i_1\\ 
i_2\\ 
\vdots\\ 
i_N
\end{pmatrix}
\]

The simulator, and devices model the non-linear current/voltage behavior
between all nodes

as such, the \(G\)'s may be non-linear functions, and include the
\(v\)'s and \(i\)'s.

Transient analysis use numerical methods to compute time evolution

The time step is adjusted automatically, often by proprietary
algorithms, to trade accuracy and simulation speed.

The numerical methods can be forward/backward Euler, or the others
listed below.

\begin{itemize}
\tightlist
\item
  \href{https://aquaulb.github.io/book_solving_pde_mooc/solving_pde_mooc/notebooks/02_TimeIntegration/02_01_EulerMethod.html}{Euler}
\item
  \href{https://aquaulb.github.io/book_solving_pde_mooc/solving_pde_mooc/notebooks/02_TimeIntegration/02_02_RungeKutta.html}{Runge-Kutta}
\item
  \href{https://en.wikipedia.org/wiki/Crank\%E2\%80\%93Nicolson_method}{Crank-Nicolson}
\item
  Gear \citeproc{ref-gear71}{{[}1{]}}
\end{itemize}

If you wish to learn more, I would recommend starting with the original
paper on analog transient analysis.

\href{https://www2.eecs.berkeley.edu/Pubs/TechRpts/1973/ERL-m-382.pdf}{SPICE
(Simulation Program with Integrated Circuit Emphasis)} published in 1973
by Nagel and Pederson

The original paper has spawned a multitude of commercial, free and open
source simulators, some are listed below.

If you have money, then buy Cadence Spectre. If you have no money, then
start with ngspice.

\textbf{Commercial} -
\href{https://www.cadence.com/ko_KR/home/tools/custom-ic-analog-rf-design/circuit-simulation/spectre-simulation-platform.html}{Cadence
Spectre} - \href{https://eda.sw.siemens.com/en-US/ic/eldo/}{Siemens
Eldo} -
\href{https://www.synopsys.com/implementation-and-signoff/ams-simulation/primesim-hspice.html}{Synopsys
HSPICE}

\textbf{Free} - \href{http://aimspice.com}{Aimspice} -
\href{https://www.analog.com/en/design-center/design-tools-and-calculators/ltspice-simulator.html}{Analog
Devices LTspice} - \href{https://xyce.sandia.gov}{xyce}

\textbf{Open Source} - \href{http://ngspice.sourceforge.net}{ngspice}


{
\centering
\aicfig{media/mixed_simulator.pdf}
}

\emph{Figure 1: Mixed signal simulator, where a digital and an analog
simulator exchange values through DAC and ADC connect modules under a
shared event and timestep control}

\section{Demo}\label{demo}

Tutorial at \url{http://analogicus.com/jnw_sv_sky130a/}

Repository at \url{https://github.com/wulffern/jnw_sv_sky130a}

Assumes knowledge of
\href{https://analogicus.com/aic2026/sky130nm_tutorial}{Tutorial}

\section{The circuit}\label{the-circuit}\index{the circuit@The circuit}

In \texttt{design/JNW\_SV\_SKY130A/JNWSW\_CM.sch} you'll find a current
mirror, and a 5-bit current DAC.

What we want from the digital is to control the binary value of the
current DAC.


{
\centering
\aicfig{media/JNWSW_CM.pdf}
}

\emph{Figure 2: The JNWSW\_CM schematic - a current mirror from ibp plus
a 5-bit binary weighted current DAC controlled by b{[}4:0{]}}

\section{The digital code}\label{the-digital-code}\index{the digital code@The digital code}

The digital code is shown below. The \texttt{clk} controls the stepping,
while the \texttt{reset} sets the output \texttt{b=0}. When
\texttt{reset} is off, then the \texttt{b} increments.

\begin{Shaded}
\begin{Highlighting}[]
\KeywordTok{module}\NormalTok{ dig}\OperatorTok{(}
           \DataTypeTok{input} \DataTypeTok{wire}\NormalTok{         clk}\OperatorTok{,}
           \DataTypeTok{input} \DataTypeTok{wire}\NormalTok{         reset}\OperatorTok{,}
           \DataTypeTok{output}\NormalTok{ logic }\OperatorTok{[}\DecValTok{4}\OperatorTok{:}\DecValTok{0}\OperatorTok{]}\NormalTok{ b}
           \OperatorTok{);}

\NormalTok{   logic                      rst }\OperatorTok{=} \DecValTok{0}\OperatorTok{;}

\NormalTok{   always\_ff }\OperatorTok{@(}\KeywordTok{posedge}\NormalTok{ clk}\OperatorTok{)} \KeywordTok{begin}
      \KeywordTok{if}\OperatorTok{(}\NormalTok{reset}\OperatorTok{)}
\NormalTok{        rst }\OperatorTok{\textless{}=} \DecValTok{1}\OperatorTok{;}
      \KeywordTok{else}
\NormalTok{        rst }\OperatorTok{\textless{}=} \DecValTok{0}\OperatorTok{;}
   \KeywordTok{end}

\NormalTok{   always\_ff }\OperatorTok{@(}\KeywordTok{posedge}\NormalTok{ clk}\OperatorTok{)} \KeywordTok{begin}
      \KeywordTok{if}\OperatorTok{(}\NormalTok{rst}\OperatorTok{)}
\NormalTok{        b }\OperatorTok{\textless{}=} \DecValTok{0}\OperatorTok{;}
      \KeywordTok{else}
\NormalTok{        b }\OperatorTok{\textless{}=}\NormalTok{ b }\OperatorTok{+} \DecValTok{1}\OperatorTok{;}
   \KeywordTok{end} \CommentTok{// dig}
\KeywordTok{endmodule}
\end{Highlighting}
\end{Shaded}

\section{Compile RTL}\label{compile-rtl}\index{compile rtl@Compile RTL}

The first thing we need to do is to translate the verilog into a
compiled object that can be used in ngspice.

\begin{Shaded}
\begin{Highlighting}[]
\BuiltInTok{cd}\NormalTok{ sim/JNWSW\_CM}
\ExtensionTok{ngspice}\NormalTok{ vlnggen ../../rtl/dig.v}
\end{Highlighting}
\end{Shaded}

\section{Import object into SPICE
file}\label{import-object-into-spice-file}

I'm lazy. So I don't want to do the same thing multiple times. As such,
I've written a small script to help me instanciate the verilog

\begin{Shaded}
\begin{Highlighting}[]
\FunctionTok{perl}\NormalTok{ ../../tech/script/gensvinst ../../rtl/dig.v dig}
\end{Highlighting}
\end{Shaded}

The script generates an \texttt{svinst.spi} file. The first section
imports the digital compiled library

\begin{Shaded}
\begin{Highlighting}[]
\ExtensionTok{adut}\NormalTok{ [clk}
\ExtensionTok{+}\NormalTok{ reset}
\ExtensionTok{+}\NormalTok{ ]}
\ExtensionTok{+}\NormalTok{ [b.4}
\ExtensionTok{+}\NormalTok{ b.3}
\ExtensionTok{+}\NormalTok{ b.2}
\ExtensionTok{+}\NormalTok{ b.1}
\ExtensionTok{+}\NormalTok{ b.0}
\ExtensionTok{+}\NormalTok{ ] null dut}
\ExtensionTok{.model}\NormalTok{ dut d\_cosim }
\ExtensionTok{+}\NormalTok{ simulation=}\StringTok{"../dig.so"}\NormalTok{ delay=10p}
\end{Highlighting}
\end{Shaded}

Turns out that ngspice needs the digital inputs and outputs to be
connected to something to calculate them (I think), so connect some
resistors

\begin{Shaded}
\begin{Highlighting}[]
\ExtensionTok{*}\NormalTok{ Inputs}
\ExtensionTok{Rsvi0}\NormalTok{ clk 0 1G}
\ExtensionTok{Rsvi1}\NormalTok{ reset 0 1G}

\ExtensionTok{*}\NormalTok{ Outputs}
\ExtensionTok{Rsvi2}\NormalTok{ b.4 0 1G}
\ExtensionTok{Rsvi3}\NormalTok{ b.3 0 1G}
\ExtensionTok{Rsvi4}\NormalTok{ b.2 0 1G}
\ExtensionTok{Rsvi5}\NormalTok{ b.1 0 1G}
\ExtensionTok{Rsvi6}\NormalTok{ b.0 0 1G}
\end{Highlighting}
\end{Shaded}

For the busses I find it easier to read the value as a real, so
translate the buses from digital \texttt{b{[}4:0{]}} to a real value
\texttt{dec\_b}

\begin{Shaded}
\begin{Highlighting}[]
\ExtensionTok{E\_STATE\_b}\NormalTok{ dec\_b 0 value=\{}\ErrorTok{(} \ExtensionTok{0} 
\ExtensionTok{+}\NormalTok{ + 16}\PreprocessorTok{*}\NormalTok{v}\ErrorTok{(}\ExtensionTok{b.4}\KeywordTok{)}\ExtensionTok{/AVDD}
\ExtensionTok{+}\NormalTok{ + 8}\PreprocessorTok{*}\NormalTok{v}\ErrorTok{(}\ExtensionTok{b.3}\KeywordTok{)}\ExtensionTok{/AVDD}
\ExtensionTok{+}\NormalTok{ + 4}\PreprocessorTok{*}\NormalTok{v}\ErrorTok{(}\ExtensionTok{b.2}\KeywordTok{)}\ExtensionTok{/AVDD}
\ExtensionTok{+}\NormalTok{ + 2}\PreprocessorTok{*}\NormalTok{v}\ErrorTok{(}\ExtensionTok{b.1}\KeywordTok{)}\ExtensionTok{/AVDD}
\ExtensionTok{+}\NormalTok{ + 1}\PreprocessorTok{*}\NormalTok{v}\ErrorTok{(}\ExtensionTok{b.0}\KeywordTok{)}\ExtensionTok{/AVDD}
\ExtensionTok{+}\KeywordTok{)}\ExtensionTok{/1000\}}
\ExtensionTok{.save}\NormalTok{ v}\ErrorTok{(}\ExtensionTok{dec\_b}\KeywordTok{)}
\end{Highlighting}
\end{Shaded}

\section{Import in testbench}\label{import-in-testbench}\index{import in testbench@Import in testbench}

An example testbench can be seen below (\texttt{sim/JNWSW\_CM/tran.spi})

\begin{Shaded}
\begin{Highlighting}[]
\ExtensionTok{...}

\ExtensionTok{.include}\NormalTok{ ../xdut.spi}
\ExtensionTok{.include}\NormalTok{ ../svinst.spi}

\ExtensionTok{*}\NormalTok{ Translate names}
\ExtensionTok{VB0}\NormalTok{ b.0 b}\OperatorTok{\textless{}}\NormalTok{0}\OperatorTok{\textgreater{}}\NormalTok{ dc 0}
\ExtensionTok{VB1}\NormalTok{ b.1 b}\OperatorTok{\textless{}}\NormalTok{1}\OperatorTok{\textgreater{}}\NormalTok{ dc 0}
\ExtensionTok{VB2}\NormalTok{ b.2 b}\OperatorTok{\textless{}}\NormalTok{2}\OperatorTok{\textgreater{}}\NormalTok{ dc 0}
\ExtensionTok{VB3}\NormalTok{ b.3 b}\OperatorTok{\textless{}}\NormalTok{3}\OperatorTok{\textgreater{}}\NormalTok{ dc 0}
\ExtensionTok{VB4}\NormalTok{ b.4 b}\OperatorTok{\textless{}}\NormalTok{4}\OperatorTok{\textgreater{}}\NormalTok{ dc 0}

\ExtensionTok{...}
\end{Highlighting}
\end{Shaded}

\section{Override default digital output
voltage}\label{override-default-digital-output-voltage}

We can override the output dac from digital to analog to ensure that the
digital signals have the right levels

\begin{Shaded}
\begin{Highlighting}[]

\ExtensionTok{*{-}}\NormalTok{ Override the default digital output bridge.}
\ExtensionTok{pre\_set}\NormalTok{ auto\_bridge\_d\_out =}
     \ExtensionTok{+} \ErrorTok{(} \StringTok{".model auto\_dac dac\_bridge(out\_low = 0.0 out\_high = 1.8)"}
     \ExtensionTok{+}   \StringTok{"auto\_bridge\%d [ \%s ] [ \%s ] auto\_dac"} \KeywordTok{)}
\end{Highlighting}
\end{Shaded}

\section{Running}\label{running}\index{running@Running}

You can run the whole thing with

\begin{Shaded}
\begin{Highlighting}[]
\BuiltInTok{cd}\NormalTok{ sim/JNWSW\_CM/}
\FunctionTok{make}\NormalTok{ typical }
\end{Highlighting}
\end{Shaded}

\section{Summary}\label{summary}

The one-page version of this chapter:

\begin{itemize}
\tightlist
\item
  SystemVerilog describes hardware as events: always\_ff becomes
  registers, combinational logic goes in always\_comb
\item
  Event-driven simulation only computes when something changes, which is
  why digital simulation is fast
\item
  Mixed-signal co-simulation bridges the two worlds: the svinst flow
  connects SPICE and the digital testbench
\end{itemize}

\section{Would you like to know
more?}\label{would-you-like-to-know-more}

The IEEE 1800 standard is the authority; Sutherland's tutorials are the
readable way in

The Verilator and iverilog manuals, for the two simulators used here

\phantomsection\label{refs}
\begin{CSLReferences}{0}{0}
\bibitem[\citeproctext]{ref-gear71}
\CSLLeftMargin{{[}1{]} }%
\CSLRightInline{C. Gear, {``Simultaneous numerical solution of
differential-algebraic equations,''} \emph{IEEE Transactions on Circuit
Theory}, vol. 18, no. 1, pp. 89--95, 1971, doi:
\href{https://doi.org/10.1109/TCT.1971.1083221}{10.1109/TCT.1971.1083221}.
}

\end{CSLReferences}
